\documentclass[12pt]{article}
\usepackage[margin=0.75in]{geometry}
\usepackage{fontspec}
\setmainfont{Latin Modern Roman}
\usepackage{microtype}
\usepackage{multicol}
\usepackage{fancyhdr}
\usepackage{titlesec}
\usepackage{enumitem}
\usepackage{xcolor}
\usepackage[hidelinks]{hyperref}
\setlength{\columnsep}{0.28in}
\setlength{\parindent}{1.1em}
\setlength{\parskip}{0.32em}
\setlength{\headheight}{15pt}
\titleformat{\section}{\large\bfseries\scshape}{}{0pt}{}
\titleformat{\subsection}{\normalsize\bfseries}{}{0pt}{}
\pagestyle{fancy}
\fancyhf{}
\fancyfoot[C]{\thepage}
\renewcommand{\headrulewidth}{0.4pt}
\newcommand{\focusbox}[1]{\par\vspace{0.4em}\noindent\begingroup\setlength{\fboxsep}{6pt}\colorbox{gray!12}{\begin{minipage}{\dimexpr\linewidth-2\fboxsep\relax}#1\end{minipage}}\endgroup\par\vspace{0.5em}}

\fancyhead[L]{Whately: Irrelevant Proof}
\fancyhead[R]{Student Reading}
\begin{document}
\begin{center}
{\LARGE\bfseries Irrelevant Proof and Missing the Point}\par\vspace{0.25em}
{\large\scshape Week 10: Answering a Question No One Asked}\par\vspace{0.5em}
{Adapted and modernized from Richard Whately, \textit{Elements of Logic}}\par\vspace{0.6em}
\rule{0.82\textwidth}{0.4pt}
\end{center}

\section*{Central Question}
How can an argument prove something and still fail to prove the point at issue?

\section*{Vocabulary Preview}
\begin{multicols}{2}
\begin{itemize}[leftmargin=*, itemsep=0.18em]
\item \textbf{Irrelevant conclusion}: a conclusion different from the one that needed proof.
\item \textbf{Ignoratio elenchi}: the traditional name for proving the wrong point or missing the issue.
\item \textbf{Point at issue}: the exact question under debate.
\item \textbf{Nearby claim}: a true or interesting claim close to the issue but not the same claim.
\item \textbf{Partial proof}: support that proves only part of what was required.
\item \textbf{Redirection}: moving attention from the disputed question to an easier one.
\item \textbf{Burden of proof}: what the speaker must establish for this argument to succeed.
\item \textbf{Prompt-match test}: asking whether the evidence answers the actual assigned question.
\item \textbf{Repair}: changing the conclusion or adding evidence so proof and claim match.
\item \textbf{Essay trap:} writing a paragraph that supports a related idea instead of the thesis.
\end{itemize}
\end{multicols}

\section*{Introduction}
Weeks 7--9 trained you to diagnose arguments that look stronger than they are: fallacy in general, equivocation, and circular reasoning. Week 10 adds a different failure. Sometimes a speaker really does prove something. The problem is that the thing proved is not the thing that was supposed to be proved.

Whately treats this as \textit{ignoratio elenchi}: ignoring the issue, or mistaking the point to be proved. It is one of the most common failures in speeches, essays, online arguments, and classroom answers.

\focusbox{\textbf{Source note:} Adapted and modernized from Whately's treatment of irrelevant conclusion and mistaken proof. Classroom examples are original. This week is especially important for student writing: many weak paragraphs prove a nearby claim instead of the assigned claim.}

\begin{multicols}{2}
\section*{Reading}

\subsection*{1. What Missing the Point Means}
An irrelevant-proof argument does not necessarily give false evidence. It may give true evidence, strong emotion, or an impressive example. But if the evidence answers a different question, the argument has missed the point.

Example:
\begin{itemize}[leftmargin=*]
\item Claim needing proof: The club should spend its budget on new debate microphones.
\item Evidence offered: The debate team works very hard.
\end{itemize}
The evidence may be true. It does not yet prove that microphones are the best use of the budget.

\subsection*{2. The Point at Issue}
The point at issue is the exact question that must be answered. Before judging an argument, state the point at issue in one sentence.

Ask: \textit{What would the speaker have to prove for the conclusion to follow?}

If the conclusion is ``we should cancel the trip,'' proving ``rain is possible'' may not be enough. The issue may be whether the rain is likely, dangerous, and not manageable by another plan.

\subsection*{3. Nearby Claims}
A nearby claim is close enough to sound relevant but different enough to leave the real question unanswered.

\begin{itemize}[leftmargin=*]
\item ``The candidate is popular'' is near ``the candidate is qualified,'' but it is not the same claim.
\item ``The book is old'' is near ``the book is wise,'' but age alone does not prove wisdom.
\item ``The rule has a good purpose'' is near ``the rule is fair in this case,'' but purpose and fairness are not identical.
\end{itemize}

\subsection*{4. Partial Proof}
Sometimes the speaker proves part of the case. That should not be ignored. But partial proof must be named honestly.

``The cafeteria needs repair'' might support spending money on the cafeteria. It does not by itself prove that every other project should be cancelled, that this contractor is best, or that the proposed cost is reasonable.

\subsection*{5. Redirection}
A speaker may redirect attention because the easier point is more attractive than the real point:
\begin{itemize}[leftmargin=*]
\item from guilt to bad character;
\item from policy cost to good intentions;
\item from whether a claim is true to whether many people believe it;
\item from whether a speech proves something to whether the speech is moving.
\end{itemize}
Redirection often works because the audience forgets the original question.

\subsection*{6. The Prompt-Match Test}
Use this four-step test:
\begin{enumerate}[leftmargin=*,itemsep=0.12em]
\item State the conclusion.
\item State the exact point at issue.
\item State what the evidence actually proves.
\item Decide whether the proof matches the conclusion; if not, repair or narrow the claim.
\end{enumerate}

\subsection*{7. How This Helps Essays}
In an essay, the point at issue is often the thesis or the paragraph's topic sentence. A paragraph can contain accurate quotations and still fail if those quotations prove a nearby point. Week 10 asks you to check whether each paragraph answers the question it promised to answer.

\subsection*{8. Literary Cases}
In \textit{Julius Caesar}, Brutus and Antony often prove true or moving points while leaving another point unsettled. Proving that Caesar was powerful is not the same as proving he had to be killed. Proving that Caesar loved Rome is not automatically the same as proving the conspirators were villains. In \textit{Macbeth}, proving that Macbeth is tempted is not the same as proving he is fated or compelled.

\subsection*{9. Repair}
There are two honest repairs:
\begin{itemize}[leftmargin=*]
\item Add evidence that answers the actual point at issue.
\item Narrow the conclusion to match what the evidence really proves.
\end{itemize}
A disciplined thinker does not pretend that a nearby proof has settled the whole question.

\subsection*{10. What This Week Adds}
A complete Week 10 diagnosis should include:
\begin{itemize}[leftmargin=*,itemsep=0.12em]
\item the conclusion;
\item the point at issue;
\item the evidence actually offered;
\item the nearby or partial claim actually proved;
\item a repaired conclusion or added evidence.
\end{itemize}

\section*{Reading Focus}
\begin{enumerate}[leftmargin=*]
\item What does \textit{ignoratio elenchi} mean in this handout?
\item How can evidence be true and still fail as proof?
\item What is the difference between the point at issue and a nearby claim?
\item Give one original example of partial proof.
\item List the four steps of the prompt-match test.
\item How does this week help student writing?
\item In \textit{Julius Caesar}, why is proving that Caesar was powerful not the same as proving he had to be killed?
\end{enumerate}

\section*{Handout Summary}
Write 5--7 sentences explaining how an argument can miss the point and how a writer should repair it. Your summary must include the words \textit{issue}, \textit{evidence}, \textit{nearby}, and \textit{repair}.
\end{multicols}
\end{document}
